\section{Transition error along the histories that a policy actually visits}
\label{exec17:s5-occupied}

The uniform envelope above has a useful virtue: it remains valid without
knowing where a transition error occurs. Its price is that a large loss range
can be charged at every possible first split. A servicing comparison can be
more informative when it preserves both the histories visited under the
policy and the financial consequences of the particular transitions that
differ. This requires additional information. It cannot be obtained by
renaming an average predictive error as a uniform transition budget.

Consider a finite horizon of $K$ transitions and a finite state $X_t$ containing
the information needed by a fixed policy. The full finite history may be used
as the state; compressing it into a balance and a delinquency band requires a
separate sufficiency assumption. Let $P_t$ be the reference transition kernel
and $Q_t$ its approximation after applying the same policy. They start from
the same distribution $\mu_0$ and use the same running rewards $r_t(x)$ and
terminal reward $g(x)$. Thus the comparison concerns model error for one
policy. It is not itself the causal effect of switching policies. Put
\begin{equation}
 \begin{split}
 v^Q_K(x)&=g(x),\\
 v^Q_t(x)&=r_t(x)+\sum_y Q_t(y\mid x)v^Q_{t+1}(y),\\
 V(Q)&=\sum_x\mu_0(x)v^Q_0(x),\\
 \mu^P_t(x)&=\Pr_P(X_t=x).
 \end{split}
 \label{exec17:s5-occupied-values}
\end{equation}
The corresponding $P$ recursion defines $v^P$ and $V(P)$. The rewards can
include active-account margin and terminal indicators for mutually exclusive
default and prepayment. The latter are sufficient for undiscounted event
losses. Discounted payments or event-time costs need the corresponding
time/state reward representation; an unchanged terminal indicator cannot
silently account for them.

Define the signed continuation residual
\begin{equation}
 b_t(x)=\sum_y\{P_t(y\mid x)-Q_t(y\mid x)\}v^Q_{t+1}(y).
 \label{exec17:s5-occupied-residual}
\end{equation}
Subtracting the two backward recursions gives
$v^P_t-v^Q_t=P_t(v^P_{t+1}-v^Q_{t+1})+b_t$. Average against $\mu^P_t$.
The first term becomes the average of the next value difference against
$\mu^P_{t+1}$. Repeating this step cancels every intermediate difference;
the terminal difference is zero because $g$ is common. Hence the exact identity
is
\begin{equation}
 V(P)-V(Q)=\sum_{t=0}^{K-1}\sum_x\mu^P_t(x)b_t(x).
 \label{exec17:s5-occupied-identity}
\end{equation}
No small-error expansion is used. The approximation's continuation value is
weighted by occupancy under the reference process. Replacing that occupancy
by the approximation's occupancy in this formula is generally incorrect.
There is a second exact identity with $Q$ occupancy and $P$ continuation
values, obtained by reversing the subtraction. It does not make the unknown
reference continuation values available from an approximate simulator.

For a row $(t,x)$, let $\epsilon_t(x)$ bound its total variation and let
$R_t(x)$ be the range of $v^Q_{t+1}$ over a set containing the support of the
signed row difference. That signed difference has total mass zero. Its
positive and negative masses are equal to its total variation, so subtracting
the smallest continuation value proves
\begin{equation}
 |b_t(x)|\leq \epsilon_t(x)R_t(x),\qquad
 |V(P)-V(Q)|\leq
 \sum_{t,x}\mu^P_t(x)\epsilon_t(x)R_t(x).
 \label{exec17:s5-occupied-envelope}
\end{equation}
Using the full next-state space gives a valid range even when the signed
support is unknown. A smaller range is valid only when the excluded
transitions are known to have zero signed discrepancy. Zero observed counts
or equal rounded probability estimates do not establish that restriction.
The financial range is a continuation-value range, rather than merely the
current period's margin. An error that changes a default probability can
alter a terminal loss even when current margin is unchanged.

An exact three-period law illustrates the distinction. Begin in active state
$A$, with absorbing states $D$ and $P$ for default and prepayment. Both laws
earn USD 20 at each active opening and assign terminal rewards
$g(A)=0$, $g(D)=-1000$ and $g(P)=-50$. Their active transition rows are
\begin{equation}
 \begin{array}{c|ccc}
       & A & D & P\\\hline
 P_t(\cdot\mid A)&0.70&0.10&0.20\\
 Q_t(\cdot\mid A)&0.70&0.11&0.19
 \end{array}
 \qquad t=0,1,2.
 \label{exec17:s5-occupied-laws}
\end{equation}
The approximation moves probability $0.01$ from prepayment to default. It
preserves active survival exactly, so both laws have opening occupancies
$(1,0.7,0.49)$ and restricted active time $2.19$. Their period-three event
vectors $(F_D,F_P,S)$ are $(0.219,0.438,0.343)$ and
$(0.2409,0.4161,0.343)$. Their values are respectively $-197.1$ and
$-217.905$ USD. A correct survival curve alone therefore misses a model-value
error of $20.805$ USD per initial account.

The continuation values make the error traceable:
\begin{center}
\begin{tabular}{c r r r r}
\toprule
$t$ & $\mu^P_t(A)$ & $v^Q_{t+1}(A)$ & $b_t(A)$ & full next-state range\\
\midrule
0 & 1 & $-169.15$ & $9.5$ & $950$\\
1 & $0.7$ & $-99.5$ & $9.5$ & $950$\\
2 & $0.49$ & $0$ & $9.5$ & $1000$\\
\bottomrule
\end{tabular}
\end{center}
At every active row the residual is
$(-0.01)(-1000)+(0.01)(-50)=9.5$ USD. Equation
\eqref{exec17:s5-occupied-identity} gives
$9.5(1+0.7+0.49)=20.805$ USD. The full-range occupied envelope is
$0.01\{950+0.7(950)+0.49(1000)\}=21.05$ USD. If the equality of the
active probabilities is part of the admitted model class, the signed error
can involve only $D$ and $P$, whose continuation range is $950$ USD at every
step. The restricted-support envelope then equals the exact $20.805$ USD.

By comparison, the earlier uniform first-split envelope, using $\epsilon=0.01$,
a running range of USD 20 and terminal range of USD 1000, is
\begin{equation}
 20(0+0.01+0.0199)+1000(0.029701)=30.299\ \mathrm{USD}.
 \label{exec17:s5-occupied-comparison}
\end{equation}
Both bounds are valid in this construction. The tighter answer uses
information about occupancy and the destination of the error that the
uniform envelope did not require. It is not a universally better bound in
every process. Complete event-history enumeration, the backward recursion
and the separate geometric-duration calculation give the same reference
values here. These are exact authored laws, without a fitted external
estimator or observed bank transition audit.~\citep{exec17_s5_occupied}

\section{Which histories can validate the transition budget?}
\label{exec17:s5-transport}

A useful error budget must refer to the histories reached under the requested
policy. A held-out prediction sample often has a different history mix.
Even when its outcomes truthfully measure the desired transition, averaging
its errors under the wrong mix can hide financially important failures.
This distribution question remains after any causal identification argument
has supplied the relevant conditional transition law.

At a fixed time, let $\rho_t(x)$ be a validation input distribution and suppose
$\mu^P_t(x)>0$ implies $\rho_t(x)>0$. For the true signed continuation
residual $b_t$, write $w_t(x)=\mu^P_t(x)/\rho_t(x)$. Cauchy--Schwarz gives
\begin{equation}
 \begin{split}
 \left|\sum_x\mu^P_t(x)b_t(x)\right|
 &=\left|\sum_x\rho_t(x)w_t(x)b_t(x)\right|\\
 &\leq C_t\left\{\sum_x\rho_t(x)b_t(x)^2\right\}^{1/2},\qquad
 C_t^2=\sum_x\frac{\mu^P_t(x)^2}{\rho_t(x)}.
 \end{split}
 \label{exec17:s5-transport-bound}
\end{equation}
Thus a validated upper bound on the squared residual can be transported only
with an appropriate occupancy factor. The inequality concerns the true
conditional residual, not an arbitrary fitted prediction loss. Estimating
that residual or its squared norm with noisy transition observations requires
its own statistical control. A small in-sample squared loss cannot simply be
inserted as the right-hand side of a guaranteed interval.

Consider two history categories with validation probabilities $(0.99,0.01)$,
policy occupancies $(0.1,0.9)$ and residuals $(0,100)$ USD. Validation's mean
absolute residual is only USD 1, while the policy's signed value residual is
USD 90. The validation squared norm is $100$ USD$^2$, and
\begin{equation}
 C_t^2=\frac{0.1^2}{0.99}+\frac{0.9^2}{0.01}
      =\frac{8020}{99},\qquad
 C_t\|b_t\|_{2,\rho}=\sqrt{802000/99}\approx90.0056\ \mathrm{USD}.
 \label{exec17:s5-transport-example}
\end{equation}
The transported bound correctly reflects the nearly complete movement of
decision-relevant occupancy into the rare validation category. It does not
interpret the small unweighted USD 1 as a policy-value guarantee. Such a
residual can be generated by changing a transition to a state whose
continuation reward differs by USD 100; the numbers describe a possible
finite transition law, not merely an algebraic pair of incompatible scores.

If validation never visits the second category, the density ratio there is
undefined. No finite value of $C_t$ repairs that lack of support. With target
occupancy $0.9$ and an otherwise unrestricted continuation range of USD 1000,
the unseen category alone can contribute up to USD 900 in absolute error.
The missing information can be addressed by sampling that category, changing
the policy to avoid it, imposing and defending a transition restriction, or
retaining the resulting broad error set. The interface should expose this
choice. Increasing draws from the first category supplies none of the
missing second-category transition information.~\citep{exec17_s5_occupied}

This support failure should be distinguished from two other problems. First,
an observational transition conditional on a chosen action need not equal
the interventional transition needed by the policy if hidden confounding
remains. Reweighting the history mix cannot resolve that scientific failure.
Second, a measured account state can become available only after the
servicing action. Conditioning on its corrected warehouse value can define
the wrong information history even with complete statistical support. A
validation design must establish transition meaning and the availability
clock before it establishes precision for a residual.

These distinctions suggest concrete uses for public and generated data. A
generated finite process can reveal all transition rows and policy
occupancies, so it can test both the exact identity and deliberate support
failures. A public longitudinal cohort can exercise history construction,
censoring, event order and reweighting diagnostics. If its assignment or
measurement process does not identify the requested causal rows, it can still
test those pipeline operations while leaving causal validation open. Public
data do not become bank transition truth through a change of column names.
Conversely, missing bank truth need not delay a useful semantic prototype
with a fully specified generated process.

\section{A finite validation interval and its longitudinal allocation}
\label{exec17:s5-validation}

The exact identity can support a prospective validation interval under
stronger, explicit sampling conditions. Fix $Q$ using data separate from
the validation sample, so its continuation values are fixed when validation
begins. Index the relevant time/state rows by $j=1,\ldots,m$. Suppose the
reference occupancies $w_j=\mu^P_{t_j}(x_j)$ are known, and obtain $n_j$
independent next-state draws $Y_{jk}$ from the specified reference row $P_j$.
The sample sizes are fixed before observing these draws. Independence
between different rows is not needed for the union bound below, but the
within-row bounded independent-draw premise is needed. Define
\begin{equation}
 \widehat b_j=\frac1{n_j}\sum_{k=1}^{n_j}v^Q_{t_j+1}(Y_{jk})
              -\sum_y Q_j(y)v^Q_{t_j+1}(y),\qquad
 \widehat B=\sum_jw_j\widehat b_j.
 \label{exec17:s5-validation-residual}
\end{equation}
Let $R_j$ bound the range of the sampled continuation value under the
reference row. The bounded-score concentration argument used earlier gives
\begin{center}
$\Pr(|\widehat b_j-b_j|>R_j\sqrt{\log(2m/\alpha)/(2n_j)})\leq\alpha/m$.
\end{center}
Apply the union bound over the rows and then the triangle inequality with
nonnegative occupancies. With probability at least $1-\alpha$,
\begin{equation}
 \left|\{V(P)-V(Q)\}-\widehat B\right|\leq
 H(\boldsymbol n):=
 \sqrt{\frac{\log(2m/\alpha)}2}
 \sum_{j=1}^m\frac{w_jR_j}{\sqrt{n_j}}.
 \label{exec17:s5-validation-interval}
\end{equation}
Under these premises $V(Q)+\widehat B\pm H$ is a finite reference-value
interval. The result is conditional on the separately trained model and the
known occupancies. It validates the model's value error from genuine
reference transitions; generating extra draws from $Q$ itself estimates a
different expectation and cannot validate $P-Q$.

Known occupancy is a material premise, not a harmless convenience. In the
three-period example it follows from the scientifically imposed active
probability $0.7$ under both laws. The validation class may allow the
default/prepayment allocation to differ by time, while retaining that known
active probability. One can then validate each active row separately without
assuming that its default fraction is constant across times. When occupancy
is estimated instead, its error multiplies the row residual. A valid analysis
needs a joint confidence region or another justified bound for these weights;
plugging estimated weights into Equation \eqref{exec17:s5-validation-interval}
does not preserve its stated guarantee automatically.

The independent draws can be obtained from a generated reference process,
using distinct trajectories or separate conditional-transition draws.
Observational subject histories generally do not provide arbitrary numbers
of independent transitions at chosen histories. Repeated rows from one
account, adaptive stopping after a favorable residual, estimated causal
weights and random survivor counts require corresponding sampling arguments.
They are not covered by declaring $n_j$ to be the number of available rows.
The public prototype can test the finite result in its declared experiment
before proposing the additional assumptions and costs of authentic
validation.

The interval also shows why equal validation effort at every time can be
inefficient. Put $a_j=w_jR_j$. If one row observation costs $c_j>0$ and the
available variable budget is $C$, consider the continuous planning problem
\begin{equation}
 \min_{n_j>0}\ \sum_j a_jn_j^{-1/2}
 \quad\text{subject to}\quad \sum_jc_jn_j=C.
 \label{exec17:s5-validation-design}
\end{equation}
For positive $a_j$ the objective is strictly convex. Its first-order
conditions are $a_j/(2n_j^{3/2})=\lambda c_j$. Solving them and using the
budget equality gives
\begin{equation}
 S=\sum_j a_j^{2/3}c_j^{1/3},\qquad
 n_j^*=\frac{C(a_j/c_j)^{2/3}}S,\qquad
 H^*=\sqrt{\frac{\log(2m/\alpha)}{2C}}\,S^{3/2}.
 \label{exec17:s5-validation-allocation}
\end{equation}
Convexity establishes the global continuous minimum. The $2/3$ exponent
arises from adding row confidence radii, each proportional to $n_j^{-1/2}$.
It need not equal an allocation that minimizes the variance of a different
estimator. Zero-occupancy or zero-range rows require no observations for this
particular value bound, although other scientific and operating checks may
still need them. Minimum row counts, capacities and integer observations
add constraints; the continuous solution is then a planning reference.

In the example, use three active rows, $\alpha=0.05$,
$w=(1,0.7,0.49)$ and full continuation ranges $(950,950,1000)$ USD. With
equal row costs and budget $30000$, equal allocation gives $10000$ draws at
each time and a simultaneous halfwidth of approximately USD $32.5680$.
The continuous allocation is approximately $(12337.92,9726.89,7935.18)$.
The feasible integer choice $(12338,9727,7935)$ uses the same budget and
has halfwidth USD $32.1752$.

For costs $(1,2,4)$ and budget $60000$ in the stipulated cost units, equal
counts of $8571$ spend $59997$ and give halfwidth USD $35.1783$. The
continuous allocation is approximately $(19905.57,9885.99,5080.61)$.
A feasible rounding and local leftover-allocation rule yields
$(19908,9886,5080)$, spends $60000$, and gives halfwidth USD $31.4016$.
This integer rule is feasible; no claim of exact integer optimality follows
from the continuous proof. The cost units describe transition-validation
effort. They do not establish staff prices, implementation costs or bank
return on investment.~\citep{exec17_s5_occupied}

These are prospective radii from fixed counts, not intervals observed in a
new validation experiment. Even if the estimated center happened to equal
the exact USD $20.805$ discrepancy, a USD $31.4016$ halfwidth would not
certify its sign. A useful design must therefore relate the financial
precision requirement to the sampling cost before commissioning expensive
validation. Better use of event-support restrictions or a justified
homogeneous-row model may reduce that cost, while introducing additional
scientific premises. The calculation makes those tradeoffs visible without
pretending that a large nominal sample guarantees a useful decision.

Finally, the interval addresses one fixed policy, a fixed continuation model
and the specified reference transition experiment. Two policy values require
joint control over their rows and their common or separate reference data.
Choosing a policy or fitting $Q$ using the same validation observations
changes the experiment. Numerical approximation of $v^Q$, reward-model
error, unsupported tails, uncertain natural-action semantics and unmeasured
assignment truth also remain distinct. They can receive justified additional
envelopes or explicit unavailable components; none is already included in
the row-sampling radius. An integration that preserves these distinctions
can explain exactly what a public prototype established and which
information an eventual bank experiment would still need.
